\documentclass[../../../include/open-logic-section]{subfiles}

\begin{document}

\olfileid[pt]{sth}{replacement}{refproofs}
\olsection{Apêndice: Resultados sobre Substituição}

A seguir, provaremos a Reflexão dentro de $\ZF$. Provaremos também
em que sentido a Reflexão é equivalente à Substituição. E
provaremos uma consequência interessante de tudo isto, relativa à
força da Reflexão/Substituição. \emph{Aviso: isto é, de longe, a parte mais avançada de matemática deste manual.} 

Começamos com um lema que, por brevidade, recorre ao expediente notacional
do \emph{sobrelinhado} para lidar com sequências de variáveis ou
objetos. Assim: ``$\overline{a}_k$'' abrevia ``$a_{k_1}$, \dots,
$a_{k_n}$'', onde $n$ é determinado pelo contexto.

\begin{lem}\ollabel{lemreflection}
Para cada $1 \leq i \leq k$, seja $\phi_i(\overline{v}_i, x)$ uma
fórmula. Então, para cada $\alpha$
existe algum $\beta > \alpha$ tal que, para quaisquer $\overline{a}_1,
\ldots, \overline{a}_k \in V_\beta$ e cada $1 \leq i \leq k$:
\[
	\exists x\phi_i(\overline{a}_i, x) \rightarrow (\exists x \in V_\beta) \phi_i(\overline{a}_i, x)
\]
\end{lem}

\begin{proof}
Definimos um termo $\mu$ como segue: $\mu(\overline{a}_1, \ldots,
\overline{a}_k)$ é o estágio mínimo, $V$, que satisfaz todos os
seguintes condicionais, para $1 \leq i \leq k$:
\[
\exists x\phi_i(\overline{a}_i, x) \rightarrow (\exists x \in V) \phi_i(\overline{a}_i, x))
\]
É fácil verificar que $\mu(\overline{a}_1, \ldots, \overline{a}_k)$ existe para todos $\overline{a}_1, \ldots, \overline{a}_k$. Ora, usando a Substituição e o nosso teorema de recursão, definimos:
\begin{align*}
	S_0 & = V_{\alpha+1}\\
	S_{n+1} & = S_n \cup \bigcup
	\Setabs{\mu(\overline{a}_1, \ldots, \overline{a}_k)}
	{\overline{a}_1, \ldots, \overline{a}_k \in S_n} \\
	S &= \bigcup_{m < \omega} S_n.
\end{align*}
Cada $S_n$, e portanto o próprio $S$, é um estágio depois de $V_\alpha$. Fixemos
$\overline{a}_1$, \dots,~$\overline{a}_k \in S$; logo existe algum $n <
\omega$ tal que $\overline{a}_1$, \dots, $\overline{a}_k \in S_n$.
Fixemos algum $1 \leq i \leq k$, e suponhamos que $\exists x
\phi_i(\overline{a}_i,x)$. Então $(\exists x \in \mu(\overline{a}_1,
\ldots, \overline{a}_k))\phi_i(\overline{a}_i, x)$ por construção, logo
$(\exists x \in S_{n+1})\phi_i(\overline{a}_i, x)$ e portanto
$(\exists x \in S)\phi_i(\overline{a}_i, x)$. Logo $S$ é o nosso $V_\beta$.
\end{proof}
\noindent 
Podemos agora provar \olref[sth][replacement][ref]{reflectionschema} de
forma bastante direta:

\begin{proof}[Prova] 
Fixemos $\alpha$. Sem perda de generalidade, podemos assumir que os únicos conectivos de $\phi$ são $\exists$, $\lnot$ e $\land$ (pois estes são
expressivamente adequados). Sejam $\psi_1, \ldots, \psi_k$ uma enumeração de
cada uma das subfórmulas de $\phi$ por ordem de complexidade, de modo que $\psi_k = \phi$.
Por \olref{lemreflection}, existe um $\beta > \alpha$ tal que, para
qualquer $\overline{a}_i \in V_\beta$ e cada $1 \leq i \leq k$:
\begin{align}\label{reflectionnicelybehaved}
	\exists x\psi_i(\overline{a}_i, x) \rightarrow 
	(\exists x \in V_\beta) \psi_i(\overline{a}_i, x)\tag{*}
\end{align}
Por indução na complexidade de $\psi_i$, mostraremos que
$\psi_i(\overline{a}_i) \leftrightarrow
\psi_i^{V_\beta}(\overline{a}_i)$, para qualquer  $\overline{a}_i \in
V_\beta$. 	Se $\psi_i$ é atómica, isto é trivial. A bicondicional
estabelece também que, quando $\psi_i$ é uma negação ou conjunção de
subfórmulas que satisfazem esta propriedade, a própria $\psi_i$ satisfaz esta
propriedade. Logo o único caso interessante é o da quantificação. Fixemos
$\overline{a}_i \in V_\beta$; então:
\begin{align*}
	(\exists x \psi_i(\overline{a}_i, x))^{V_\beta}
	&\text{ sse }
	(\exists x \in V_\beta)\psi_i^{V_\beta}(\overline{a}_i, x)
	&&\text{por definição}\\
	&\text{ sse }
	(\exists x \in V_\beta)\psi_i(\overline{a}_i,  x)
	&&\text{por hipótese}\\
	&\text{ sse }
	\exists x \psi_i(\overline{a}_i, x)
	&&\text{por \eqref{reflectionnicelybehaved}}
\end{align*}
Isto completa a indução; o resultado segue pois $\psi_k = \phi$.
\end{proof}

Provamos a Reflexão em $\ZF$. A nossa prova seguiu essencialmente
\citet{Montague1961}. Queremos agora provar em $\Z$ que a
Reflexão implica a Substituição. A prova segue \citet{Levy1960},
com uma simplificação. 

Como trabalhamos em $\Z$, não podemos apresentar a Reflexão exatamente na
forma dada acima. Afinal, formulamos a Reflexão usando a
notação ``$V_\alpha$'', e isso não pode ser definido em $\Z$ (ver
\olref[sth][spine][zf]{sec}). Em vez disso, ofereceremos uma formulação
aparentemente mais fraca da Substituição, como segue:

\begin{defish}
\emph{Reflexão fraca.} Para qualquer fórmula $\phi$, existe um conjunto
transitivo $S$ tal que $0$, $1$ e quaisquer parâmetros de $\phi$ são
!!{element}os de $S$, e $(\forall \overline{x} \in S)(\phi \liff
\phi^S)$.
\end{defish}

Para usar isto na prova da Substituição, seguiremos primeiro \citet[primeira parte do Teorema 2]{Levy1960} e mostraremos que podemos ``refletir'' duas
fórmulas de uma vez:

\begin{lem}[em $\Z + \text{Reflexão fraca}$.]\ollabel{lem:reflect}
Para quaisquer fórmulas $\psi, \chi$, existe um conjunto transitivo $S$ tal que
$0$ e $1$ (e quaisquer parâmetros das fórmulas) são !!{element}os de
$S$, e $(\forall \overline{x} \in S)((\psi \liff \psi^S) \land (\chi
\liff \chi^S))$.
\end{lem}

\begin{proof}
Seja $\phi$ a fórmula $(z = 0 \land \psi) \lor (z = 1 \land \chi)$. 

Aqui usamos uma abreviatura; devemos explicitar ``$z = 0$'' como
``$\forall t\, t \notin z$'' e ``$z =1$'' como ``$\forall s(s \in z
\liff \forall t\, t \notin s)$''. Mas como $0, 1 \in S$ e $S$ é
transitivo, estas fórmulas são \emph{absolutas} para $S$; isto é, aplicam-se ao mesmo objeto quer restrinjamos os seus quantificadores a
$S$.\footnote{Mais formalmente, seja $\xi$ qualquer uma destas
fórmulas; então $\xi(z) \liff \xi^S(z)$.}

Pela Reflexão fraca, temos algum $S$ adequado tal que:
\begin{align*}
	(\forall z, \overline{x} \in S)(&\phi \liff \phi^S)\\
	\text{O.s.\ }(\forall z, \overline{x} \in S)(&((z = 0 \land \psi) \lor (z = 1 \land \chi)) \liff {}\\
	&\phantom{(}((z = 0 \land \psi) \lor (z = 1 \land \chi))^S)\\
	\text{O.s.\ }(\forall z, \overline{x} \in S)(&((z = 0 \land \psi) \lor (z = 1 \land \chi))\liff {}\\
	&\phantom{(}((z = 0 \land \psi^S) \lor (z = 1 \land \chi^S)))\\
	\text{O.s.\ }(\forall \overline{x} \in S)(&(\psi \liff \psi^S) \land (\chi \liff \chi^S))
\end{align*}
A segunda afirmação implica a terceira porque ``$z = 0$'' e ``$z=1$''
são absolutas para $S$; a quarta segue pois $0 \neq 1$.
\end{proof}\noindent Podemos agora obter a Substituição, seguindo e simplificando 
\citet[Teorema 6]{Levy1960}:

\begin{thm}[em $\Z$ + Reflexão fraca]\label{thm:replacement} 
Para qualquer fórmula $\phi(v,w)$, e qualquer $A$, se $(\forall x \in A)\lexists![y][\phi(x,y)]$, então
$\Setabs{y}{(\exists x \in A)\phi(x,y}$ existe.
\end{thm}

\begin{proof}
Fixemos $A$ tal que $(\forall x \in A)\lexists![y][\phi(x,y)]$, e
definimos fórmulas:
\begin{align*}
	\psi &\text{ é } (\phi(x, z) \land A = A)\\
	\chi &\text{ é } \lexists[y][\phi(x, y)]
\end{align*}
Usando \olref{lem:reflect}, como $A$ é um parâmetro de $\psi$, existe
um conjunto transitivo~$S$ tal que $0, 1, A \in S$  (juntamente com quaisquer outros
parâmetros), e tal que:
\[
	(\forall x,z \in S)((\psi \liff \psi^S) \land (\chi \liff \chi^S))
\]
Logo, em particular:
\begin{align*}
	(\forall  x, z \in S)(&\phi(x, z) \liff \phi^S(x, z))\\
	(\forall x \in S)(&\exists y\phi(x, y) \liff (\exists y \in S)\phi^S(x, y)) 
\end{align*}
Combinando isto, e observando que $A \subseteq S$ pois $A \in S$ e $S$ é transitivo:
\begin{align*}
	(\forall x \in A)(&\exists y\phi(x, y) \liff (\exists y \in S)\phi(x, y))
\end{align*}
Ora, $(\forall x \in A)(\lexists![y \in S])\phi(x, y)$, porque
$(\forall x \in A)\lexists![y][\phi(x, y)]$. Ora, a Separação dá
$\Setabs{y \in S}{(\exists x \in A) \phi(x, y)} = \Setabs{y}{(\exists
x \in A) \phi(x, y)}$. 
\end{proof}

\end{document}
